\subsection{Descent of the row lattice and its maximal slope}
\label{subsec:fitted-row-height}

The matrix estimates refer to a different tuple at each assigned place. Descent assembles their integral images into the base row on which the maximal slope is taken. Let \(K\) be its
totally complex coefficient field, \(t_0\) the base tuple,
and \(t(w)\) the tuple assigned to a finite place \(w\).
Write \(P_t:V_t\twoheadrightarrow Q_t\) for the column quotient,
\(\iota_t:Q_t\to R_{0,K}\) for the row comparison, and \(M_t\) for
the integral lattice in \(V_t\). The base comparison \(\iota_{t_0}\)
is an isomorphism. The selected row has
\[
 (R_0)_w=\iota_{t(w)}P_{t(w)}(M_{t(w),w}).
\]
Here and below a subscript \(w\) on a lattice means localization.

The coefficient tests of Section~\ref{subsec:fitted-evaluation} give the matrix \(V(t)\) over \(K\) on these tuple coordinates.
Write \(r_t=\prod_{ij}r_{t,ij}\) for the
product of the square ratios in the field fit. The base ratios
satisfy \(r_{t_0,ij}=1\), and the ratios at realized tuples are
nonzero.

\begin{proposition}[Height of the descended base row]
\label{prop:descended-row-height}
Use the nonvanishing in Lemma~\ref{lem:fitted-completed-nonvanishing} at the selected point and the evaluation grid and Fourier fit of the realization supplying the base quotient and integral frame. For this row estimate one may instead fix admissible tests on the same grid with the uniform size estimates of Proposition~\ref{prop:fitted-grid-constituent-costs}, the common tuple comparisons and both tests of Lemma~\ref{lem:fitted-base-minor} at the base. At every selected point above one threshold, the descended base row satisfies
\[
 \mu_{\max}(R_0)+\bar\kappa\le b x^{C_H},
\]
where
\(\bar\kappa\) is the mean logarithm of its period scalar.
Here \(C_H\ge16(\lceil|C_t|\rceil+10)+1\) suffices. The sharp
level passage adds at most \(14\log|G|\), and the resulting
bound is at most \(b x^{C_H+1}\) when \(x\ge56\).
\end{proposition}

\begin{proof}
\proofstep{Descend one comparison map on all realized tuples.}
At the selected place let \(j_t\) be the complex realization of
\(Q_t\) in the common theta space. The primitive links identify
the images of the base and tuple realizations. The row comparison
is the relative matrix for this inclusion, expressed over \(K\)
in the chosen row coordinates. In particular \(j_{t_0}\) is
injective and, for each \(u\in Q_t\),
\[
    j_{t_0}\bigl(\sigma_0(\iota_t u)\bigr)
                  =j_t\bigl(\sigma_0(u)\bigr).
\]
Here \(\sigma_0\) denotes coordinate embedding into the complex
realization. The coefficient test uses the same vector \(\mathcal F\)
of convergent level functionals on that common space. Thus its
reading on \(j_t(\sigma_0(P_t y))\) equals its reading on
\(j_{t_0}(\sigma_0(\iota_t P_t y))\). Write \(F\) for the latter
test as a function on the embedded base row. In the matrix
coordinates of the test, the integral comparison is

\[
  z_t c_t F(\iota_t P_t y)=c_* V(t)y.
\]
The \(z_t\) are roots of unity in \(K\), each \(c_t\) is nonzero,
and \(c_t/c_{t_0}\) is the image of \(r_t\).
Here the displayed expression uses the embedded row coordinates,
including the conjugate embedding in the theta test. The
semilinearity of the test with respect to that embedding gives
ordinary complex matrix multiplication in these coordinates.
The base quotient has positive rank and its nonzero vectors are
separated by the level functionals. Surjectivity of \(P_{t_0}\)
therefore implies \(c_*\ne0\). The same comparison at the base
shows that \(\ker V(t_0)=\ker P_{t_0}\).

Put \(V'(t)=(z_t/z_{t_0})^{-1}V(t)\).
Multiplying the level functionals by \(z_{t_0}c_{t_0}/c_*\)
gives a single linear map at the selected embedding. Its
matrix \(\Psi\) satisfies
\begin{equation}\label{eq:fitted-common-fourier}
 \Psi\iota_t P_t=r_t^{-1}V'(t).
\end{equation}
Choose a \(K\)-linear right inverse \(A_0\) to
\(\iota_{t_0}P_{t_0}\). Since \(r_{t_0}=1\) and
\(V'(t_0)=V(t_0)\), the base equation reads
\[
                   \Psi=V(t_0)A_0.
\]
It determines the matrix over \(K\), and the injective complex
embedding descends every equation \eqref{eq:fitted-common-fourier}
to \(K\).

If \(\Psi x=0\), choose \(y\) with
\(\iota_{t_0}P_{t_0}y=x\). Then \(V(t_0)y=0\), so
\(P_{t_0}y=0\) and \(x=0\).
Thus the kernel test factors through the two operators as
\[
  \Psi x=0\ \Longrightarrow\ V(t_0)y=0\ \Longrightarrow\ P_{t_0}y=0\ \Longrightarrow\ x=0 .
\]
 At any finite place multiplication
of \(V(t)\) by the phase ratio changes it by a unit. Thus
one can use the base archimedean bound and the local bounds on
the unnormalized matrices.

\proofstep{Assemble the images in a fractional ideal times the test dual.}
The coefficient indices group into a datum, a class
representative, and a coefficient basis in that block.
The test lattice for the original evaluation is the direct
sum of the class lattices in those coordinates. The target
\(M\) is its bilinear dual. At infinity its Hermitian Gram
is the transpose of the inverse of the test Gram. The
transpose implements bilinear evaluation. In particular,
bounding coordinates separately in arbitrary bases could
lose the dimension in the operator bound.

For the finite calculation, write \(\Sigma\) for the sum of
level depths times the logarithms of the norms of their
moduli, and set
\[
 b_1(C)=CfmN^2b(\log b+\Sigma).
\]
The factors in each expanded class term are the external
scalar, the product of local coefficients, the local datum
values, the inverse stabilizers, the two chains of length
\(m\), the initial row form, and the final pairing.
Use the local lower valuation bounds for these factors at
comparison constant \(C/16\). The first four cost functions
have weighted sum at most \(m b_1(C/16)\) each, and the
remaining four at most \(b_1(C/16)\) each before iteration.

In particular the first four contributions together are at most
\[
  \sum_{i=1}^4 m\,b_1(C/16)=4m\,b_1(C/16).
\]
The initial top test is integral. Iterating the top and row
maps multiplies their two costs by \(m\). Each term is
therefore bounded below by the negative of the sum of the
eight costs with these multiplicities. The expansions over
cosets, classes and datum support are finite and have that
same lower bound, by the ultrametric inequality.

For the finite set of class and coset terms \(a_i\), the bound passes to the reading by
\[
  v_w\left(\sum_i a_i\right)\ge\min_i v_w(a_i).
\]
 Their
cardinalities take no part in the finite cost. Pairing with
an integral vector of the test lattice and expanding it
block by block gives the same bound on every dual pairing.
Testing on a local basis gives membership in the scaled
dual lattice. Equivalently, if \(d_w^*\) is the sum of the
eight local costs, one clears denominators with a scalar
of valuation at most \(\lfloor d_w^*\rfloor\). A scalar of
that valuation times a local unit clears all other
denominators, by finite generation. The full sum of
\(d_w^*\) with finite weights is at most
\(8m b_1(C/16)=m b_1(C)/2\).

At a complex place use
\[
 a_1(C)=CfNb\log b+CmN^2\Sigma,\qquad U=Cm^2N^2\Sigma.
\]
In this paragraph the individual factor bounds are at \(C\).
Within each test block the three sums over cosets, classes
and coset-expansion support cost \(e^{ma_1(C)}\) each in
the absolute norm, and the external scalar costs one
further such factor. Inverse stabilizers are reciprocals
of positive integers and cost at most one. The coefficient
and datum bound costs \(e^U\). Let the two chain gains in
squared norm be \(e^{g_T(v)}\) and \(e^{g_R(v)}\) per step,
with both mean exponents at most \(a_1(C)\). The terminal
pairing contributes \(e^{a_1(C)}\) in squared norm. Thus for
a block of test norm squared \(B_c\) and an initial row
bound \(B_f\), the squared value is at most
\[
 e^{8ma_1(C)+2U+a_1(C)+m g_T(v)+m g_R(v)}B_cB_f.
\]
There are at most \(e^{2ma_1(C)}\) test blocks, by the
counts of data and class representatives.

Cauchy--Schwarz
uses this count once in squared norm and replaces the sum
of \(B_c\)'s by the full squared norm in the test Gram.
The initial row form contributes one more factor
\(e^{ma_1(C)}\) in \(B_f\). Taking the dual norm of the tested operator \(A_v\) at a complex place and collecting these exponents gives
\[
  \log\|A_v\|_{\rm op}^2\le (8+2+1)m a_1(C)+2U+a_1(C)+m g_T(v)+m g_R(v).
\]
Complex density of \(K\) extends the tested inequality
from algebraic vectors to the entire complex spaces.
Taking the mean and using \(U\le ma_1(C)\) gives
\(16m a_1(C)\). Increasing the comparison constant by
the factor \(16\) absorbs this overhead. The sum \(\Sigma\)
is at most the enlarged modulus budget \(S_b\) below, so
\[
 A_{\rm op}\le Cf mNb\log b+Cm^2N^2S_b.
\]
These norm estimates use the same test Gram and integral
test basis as the equation for \(\Psi\).
\label{cor:fitted-row-height-1}

\proofstep{Use the operator norm at infinity and ideal valuations at finite places.}
Let \(M\) be the coefficient lattice of the functional with its
fitted Hermitian metric. Use the square-norm slope convention, with
finite weights 
\[
  \omega_w=\log\mathrm Nw/[K:\Q]
\]
 and the uniform
mean over the complex places. The away condition in the field
fit applies to \(t(w)\) at \(w\). To see this, a prime of the
host that occurs in the tuple level is ramified for its
level character. Off the twisted places the character at
\(t(w)=t_0\) is trivial.

At a twisted place over \(\ell\)
the character is \(\delta_\ell\psi_w\). The second factor
is unramified above \(\ell\), and the first uses the deletion
prime different from \(\ell\), as in
Section~\ref{sec:transport}.
Consequently, for the inertia group \(I\) at any place above \(\ell\), the local character has
\[
  (\delta_\ell\psi_w)|_I=(\delta_\ell|_I)(\psi_w|_I)=1.
\]
 Thus every prime of this
extra level is away from \(\ell\). The finite estimate supplies
a real cost \(d_w^*\), with finite support, such that for each
integral input \(y\) there is a nonzero scalar \(s\) with
\(v_w(s)\le d_w^*\) and \(s V(t(w))y\in M\).

In the local ring, \(s\) is a unit times a power of a
uniformizer. Thus 
\[
  V(t(w))M_{t(w),w}\subset w^{-d_w}M_w
\]
with \(d_w=\lfloor d_w^*\rfloor\). The choice of \(s\) may
vary with \(y\). The floor works for every integral input
and preserves the upper bound for the weighted sum of the
costs. The product formula is used with integral valuations
\(v_w\). Equation~\eqref{eq:fitted-common-fourier}
and localization at every \(w\) give
\[
 \Psi R_0\subset JM,\qquad
 J=\prod_w w^{-v_w(r_{t(w)})-d_w}.
\]
The exponent is zero outside the twisted places and the support
of \(d_w\). It is a signed exponent. A negative contribution
may decrease the height.

We identify the metric on the parent row. Write \(v_*\) for the
infinite place of the chosen realization, \(n_P\) for its level
index and \(J_t\) for the input frame of the theta map \(\Theta_t\).
The image of \(J_t\) lies in
\(\operatorname{im}\Theta_t^*=(\ker\Theta_t)^\perp\). The projected
level and the adjoint preserve the two marked spaces. Thus the
frame reads the quotient by the kernel in its orthogonal
representatives, with
\[
 \langle J_t x,J_t y\rangle=n_P\langle x,y\rangle_{Q_t,v_*}.
\]
The relative matrix \(\iota_t\) was chosen on
\(j_t=\Theta_t J_t\). The Gram of the row at \(v_*\) is the
Gram of \(j_{t_0}\) divided by \(n_P\).

If \(j_t x\) is of
type \(p\), the coupled theta-adjoint period identity on the
selected level gives
\(\Theta_t\Theta_t^*j_t x=c_p j_t x\), where
\(c_p=\kappa(v_*)\prod_j\lambda_{j,t,p}(v_*)>0\).
The difference \(\Theta_t^*\Theta_t J_t x-c_p J_t x\) is
in the kernel of \(\Theta_t\) and in its orthogonal complement.
Consequently

\[
  \langle j_t x,j_t y\rangle=c_p n_P\langle x,y\rangle_{Q_t,v_*}
\]
for every \(y\). Dividing by \(n_P\) gives exactly the scalar
for \(\iota_t\) in the parent row. The conjugate complex
realization has the conjugate adjoint; a unit phase in the
theta map cancels against its conjugate. The individual and
coupled period identities used here concern the same projected
class sums in Sections~\ref{subsec:selected-period}
and~\ref{subsec:selected-lift-scalar}.

At another infinite place \(v\) the parent metric is defined
at the base tuple by the orthogonal sum
\[
 \langle x,y\rangle_{R_0,v}=
 \kappa(v)\sum_p\left(\prod_j\lambda_{j,t_0,p}(v)\right)
       \langle E_p x,E_p y\rangle_{Q_{t_0},v}
\]
in base coordinates. Here \(E_p\) are the rational projections
in the direct sum of the base type spaces. The base relative
matrix is the identity. The rational quotient pairing has
orthogonal distinct Hecke parts at these places by the adjoint
relations for the conjugate Hecke operators. The commuting
operators and their adjoints descend to the quotient by the
stable kernel, represented by its orthogonal complement. They
are simultaneously diagonalizable there. Distinct type parts
are sums over disjoint sets of joint eigencharacters.

For each
pair of characters in different sets choose an operator with
different eigenvalues, and apply its adjoint relation. This
gives the claimed orthogonality by sesquilinearity. Thus 
\[
  \sum_p E_p=1
\]
 recovers the full
quotient norm from these terms. The row estimate uses
the base tuple at every infinite place and the other used
tuples at \(v_*\). These are precisely the two domains of
the metric comparison. \label{cor:fitted-row-height-2}

Divide the squared metric by \(\kappa(v)\) and call the result
\(R_0^{\rm n}\). At \(v_*\) the type spaces span and are
eigenspaces of the positive comparison of the two metrics.
At the other places use the displayed sum. Denote the quotient metric by \(\|\cdot\|_{Q,v}^2\). On either
domain the comparison reads
\[
 \|z\|_{R_0^{\rm n},v}^2\ge e^{g(v)}\|z\|_{Q,v}^2
 \quad\text{if}\quad
 g(v)\le\sum_j\log\lambda_{j,t_0,p}(v)
       \quad\text{for every occurring }p .
\]

The quotient metric is the one induced by the surjection
\(P_{t_0}\). Consequently the operator norm of
\(\Psi\iota_{t_0}\) from \(Q_{t_0}\) equals the norm of its
composition with \(P_{t_0}\), namely \(V(t_0)\). If
\(A_{\rm op}\) bounds the mean logarithm of this squared norm,
the logarithmic squared operator height of \(\Psi\) from
\(R_0^{\rm n}\) is at most \(A_{\rm op}-\operatorname{mean}g\).
For the same fitted data, the bound from the individual
lift scalars is 
\[
  -\operatorname{mean}g\le
3CfmN^2\log b.
\]
 Indeed one can take the sum of the lower
logarithmic bounds for the \(m\) blocks. The hook terms are
nonnegative, and the loss per block is at most
\(3CfN^2\log b\) by the parameter bound for \(L_p\), with
\(n_V\le N\).

For a rank-\(s\) sublattice \(L\subset R_0^{\rm n}\),
the image \(\Psi L\subset JM\) has the same rank. On its
determinant line 
\[
  \|\wedge^s\Psi\|^2_v\le\|\Psi\|_v^{2s}.
\]
Since \(\Psi:L\to\Psi L\) is an isomorphism of the
underlying integral modules, the archimedean inequality gives
\[
 \mu(L)\le \mu(\Psi L)
                +\operatorname{mean}_v\log\|\Psi\|_v^2.
\]
The first term is at most \(\mu_{\max}(M)-2\log\mathrm N(J)/[K:\Q]\),
and the second is at most \(A_{\rm op}-\operatorname{mean}g\).
Taking the supremum over \(L\) gives
\begin{equation}\label{eq:fitted-row-slope}
 \mu_{\max}(R_0)+\bar\kappa
 \le \mu_{\max}(M)+A_{\rm op}-\operatorname{mean}g+
   2\sum_w\omega_w\bigl(v_w(r_{t(w)})+d_w\bigr).
\end{equation}
The sign on \(\bar\kappa=\operatorname{mean}\log\kappa\)
follows from the division of the squared metric. The factor
two in the finite contribution is from the same squared-norm
convention.

\proofstep{Apply the product formula to the fitted square ratios.}
Use \(N=k+m(t+k+m+6)\), \(r=k+m-1\), and \(B_T\)
from Section~\ref{subsec:fitted-square-ratios} for the
weighted absolute sums of the shifted toric valuations.
In each square, the two unshifted
bottom values \(A_{ij},A'_{ij}\) are nonzero. Their
valuations are at least \(-3\delta_w\), with
\(\sum_w\omega_w\delta_w\le D_b\), and their positive
archimedean logarithms have mean at most \(A_b\).

For
either value \(z\), the product formula gives
\[
 \sum_w\omega_w |v_w(z)|
 =\operatorname{mean}_v\log|z|_v+
   2\sum_w\omega_w\max(0,-v_w(z))
 \le A_b+6D_b.
\]
Set \(B=\max(A_b+6D_b,B_T)\). The two bottom terms and
four copies of each of two toric terms contribute at most
\(2(2+8N)B\) per square, allowing for the two local branches
of the toric label. The upper bound \(N\) is for the
half rank of the square. If \(B_N\) bounds the prime sum
of the density term for each of its indices, the density
part costs at most \(10NB_N\) per square. In this sum a
weight at a place over \(p\) includes \(e_w=v_w(p)\), so
the full fibre weighs precisely \(\log p\).
After summing each fibre, the contribution of the error with denominator \(p^2\) is bounded by
\[
  2N\sum_p\frac{\log p}{p^2}\le2N\sum_{n\ge2}\frac{\log n}{n^2}<2N.
\]

Indeed \((\log t)/t^2\) decreases for \(t\ge2\), so
\[
 \sum_{n\ge2}\frac{\log n}{n^2}\le
 \frac{\log2}{4}+\frac{\log3}{9}
   +\int_3^\infty\frac{\log t}{t^2}\,dt
 =\frac{\log2}{4}+\frac{\log3}{9}+\frac{1+\log3}{3}<1.
\] The defect part costs at most
\(6D_b\). The fixed prime-power moduli contribute at most
\(4mS_b\), where \(S_b\) is the sum of the prescribed
exponents times the logarithms of the moduli. Here the
norm of an integer bounds its contribution on any subset
of places. All together, the square estimate gives
\begin{equation}\label{eq:fitted-grid-sum}
 \sum_w\omega_w v_w(r_{t(w)})
 \le m^2\bigl(2(2+8N)B+10NB_N+2N+6D_b+4mS_b\bigr).
\end{equation}
The signed sum on the left requires the actual tuple
assignment at each finite place.

\proofstep{Absorb the uniform estimates in one exponent before changing level.}
Set \(L=\log b\). Define
\[
 D_*=b(k+m)g_0g_1p_dp'_d|\operatorname{disc}E|,
 \qquad P=C f^2(k+m)^6\log D_*.
\]
The prime-power budget satisfies \(0\le S_b\le2P\).
For a square \((i,j)\), let \(H_{ij}\) be its half rank and
write \(n_{ij,s}(\ell)\) for the denominator exponent chosen in
its toric comparison at \(1\le s\le H_{ij}\). Every rational prime at which
one of the selected exponents is nonzero belongs to the grid's
denominator prime set \(S_d\). The square inequality uses
\({\bf1}_{\ell\in S_d}e_w n_{ij,s}(\ell)\) at \(w\mid\ell\).
In particular at \emph{each} \((i,j,s)\) the sum to be bounded is
\[
 \sum_{\ell\in S_d}n_{ij,s}(\ell)
        \sum_{w\mid\ell}\omega_w e_w
   =\sum_{\ell\in S_d} n_{ij,s}(\ell)\log\ell .
\]
The denominator calculation of Section~\ref{subsec:fitted-square-ratios}
applies to the function supplied with this very critical integer
and base conductor. Its values on \(S_d\) are nonnegative, so
restricting to the actual twisted fibres only decreases the sum.
This supplies \(B_N\) before the sum over \(s\).

On the fitted grid choose the common bottom parameters and for
each \((i,j)\) its two definite pairs \(P_{ij},P'_{ij}\) at the
base conductor. The equations used to form the square include
\(\sigma_*(A_{ij})=\operatorname{norm}(P_{ij})\) and
\(\sigma_*(A'_{ij})=\operatorname{norm}(P'_{ij})\), where
\(\operatorname{norm}\) is the mass-normalized form value of
the pair. The bottom smallness and nonvanishing assertions
hold for each pair satisfying these bottom equations. The
finite integral comparison applied to this list gives one
nonnegative function \(\delta_w\) whose weighted sum on
every finite set is at most \(D_b\), and

\[
  v_w(A_{ij}),v_w(A'_{ij})\ge-3\delta_w.
\]
 The common
defect in \eqref{eq:fitted-square-bottom} is this same function.
For the infinite-place comparison take the conjugate norm sums
of the two selected pairs. Since these values are the images of
\(A_{ij},A'_{ij}\), averaging the positive logarithms bounds
the two field heights appearing in the product formula above.
One may compute in an extension containing the class coordinates, since
restriction of embeddings gives the same mean. The rank,
weight and base conductor in both calculations are those
recorded on this grid.
\label{cor:fitted-row-height-3}

For the two bottom norms put \(L_0=\log Q_0\), where
\(Q_0\) is the product of the base place norms with their
conductor exponents, and let \(B_w\) be the upper bound on the
fitted square weights. The bottom comparison uses the mass
sum of the same small spherical pairs as in
Section~\ref{subsec:fitted-square-ratios}. Its denominator
and conjugate size estimates give
\begin{align*}
 D_b &\le C f(2N)\log(2fN)
    +C f(2N)B_w L+C f(2N)^2 B_w L_0,\\
 A_b &\le C f(2N)B_w\bigl(L+\log(2NB_w)\bigr)
    +C fN^2\bigl(L+\log(2fNB_w)\bigr)+C N^2 L_0 .
\end{align*}
The denominator estimate uses integral representatives
at each finite place. The real estimate uses the conjugate
small spherical pairs and the mass normalization on the base
level, with the root discriminant and the host logarithmic
derivative bounded by the size record. These are estimates
on the unshifted values \(A,A'\).
\label{cor:fitted-row-height-4}
Combining them with the toric bounds gives

\[
  B_T,B_N\le Cf b(L+\log x),
\]
\(D_b\le186C^2 f^2N^2 bL\), and
\(A_b\le57C^3 fN^2 bL\). The last two estimates use
\(L_0\le45fL\) and \(1\le B_w\le Cb\) with
\(\log B_w\le CL\). In the defect bound, the three terms
are bounded respectively by \(4,2,180\) times
\(C^2 f^2N^2 bL\). For the positive real logarithm bound,
the three multipliers of \(C^3 fN^2 bL\) are \(8,4,45\).
These estimates use \(L\ge1\),
\(\log(2fN),\log(2N)\le 1+L\), and
\(\log x\le2L\).

For clarity, the two monomials that absorb all the
costs are
\[
 X_0=C^3 f m^2rN^3 bL,\qquad Y_0=C^3 f m^2N^2 bP .
\]
The coefficient lattice estimate \(4f n_*R_*\log(n_*R_*)\)
from Proposition~\ref{prop:fitted-coefficient-slope} is at most \(32X_0\) by
\(n_*\le2mr\), \(R_*\le CmN^3 b\) and
\(\log(n_*R_*)\le4L\).
The bound for the archimedean operator height, including
\(-\operatorname{mean}g\), is
\[
 Cf mN bL + C m^2N^2 S_b + 3Cf mN^2 L\le4X_0+2Y_0.
\]
The bound for twice the weighted sum of \(d_w\) is
\[
 Cf bNm^2(L+S_b)+Cf m^2N^2b(L+S_b)
 \le 2X_0+4Y_0.
\]

Since \(B\le1176C^3 f^2N^2bL\), the right side of
\eqref{eq:fitted-grid-sum} is at most
\((23520+30+2+1116)X_0+8Y_0\).
Here \(f\le r\), \(m\le N\), and all size parameters
multiplying these monomials are at least one.
Substitution into \eqref{eq:fitted-row-slope} gives the bound
\(100000(X_0+Y_0)\). This is \(b\) times the polynomial budget
in the descent. Put \(T=\lceil|C_t|\rceil+1\) and

\[
  Q=x^{T+9}.
\]
 The selected size bounds give
\(f,m,k,t,L\le Q\) and \(\log D_*\le20Q\).
In particular \(N\le10Q^2\), \(r\le2Q\), and the expression
\(rNL+P\) is at most \(20Q^4+1280CQ^9\), hence at
most \(1300CQ^9\). Multiplying by \(100000 C^3 f m^2N^2\)
bounds the polynomial by \(A C^4x^{16(T+9)}\) with
\(A=13\cdot10^9\). With \(x\ge AC^4\), one more power of \(x\) absorbs the leading coefficient:
\[
  AC^4x^{16(T+9)}\le x^{16(T+9)+1}\le x^{C_H}\qquad(C_H\ge16(T+9)+1).
\]

The fitted-field cutoff precedes this descent. It is the
maximum of the cutoff for the evaluation grid and the
cutoff for the Fourier fit. The same realization of the row
gives the base quotient, the field fit and the
integral frame. For any fixed family and fixed cutoff,
first take the maximum required comparison constant and
then the maximum eventual scale thresholds.
With the cutoff and family fixed, if \(C_i\) and \(x_i(C)\) are the comparison parameters of the two estimates, one may take
\[
  C=\max(C_1,C_2),\qquad x_0=\max\{x_1(C),x_2(C)\}.
\]
 These thresholds
are uniform over the point indices and their realized tuples.

The fitted matrix in Section~\ref{subsec:fitted-evaluation}
and the all-powers estimate in
\eqref{eq:fitted-toric-all-powers} apply to the row at this same point.
The resulting height bound is
\(\mu_{\max}(R_0)+\bar\kappa\le b x^{C_H}\)
on the parent base lattice. Its period and metric are the ones
used to define the height of the descended map. The sharp passage adds
\(14\log|G|\), which is at most \(56b\) under the group
size estimate. For \(C_H\ge1\) and \(x\ge56\) this
additional cost fits within \(x^{C_H+1}b\).
\label{cor:fitted-row-height-5}

\end{proof}

The signed covolume, the two local indices and the marked length
comparison in the window argument refer to these same selected rows.
